\documentclass[14pt,russian]{extarticle}
\usepackage[T1,T2A]{fontenc}
\usepackage[utf8]{inputenc}
\usepackage[14pt]{extsizes}
\usepackage[]{geometry}
\geometry{a4paper,
total={170mm,257mm},left=2cm,right=2cm,
top=2cm,bottom=2cm}

\usepackage[utf8]{inputenc}
\usepackage[russian]{babel}
\usepackage{amsmath}
\usepackage{hyperref}
\usepackage{graphicx}
\usepackage{amssymb}

\title{Силы трения}

\begin{document}
\maketitle

Очень полезная статья SIGGRAPH 21: \url{https://siggraphcontact.github.io/assets/files/SIGGRAPH21_friction_contact_notes.pdf}

\paragraph{Линейная задача о дополнительности}
Формулируется так:

\[
\begin{matrix}
    Ax + b \ge 0 \\
    x \ge 0 \\
    x^T (Ax + b) = 0 \\
\end{matrix}
\]

С помощью этих уравнений и неравенств удобно задать условия непроникновения:

\begin{itemize}
    \item Относительная скорость точки больше или равна нулю (тела не проникают друг в друга)
    \item Сила прикладывается неотрицательная (либо не прикладывается вообще)
    \item Если сила прикладывается, то относительная скорость должна быть равна нулю.
        Если относительная скорость больше нуля, то сила прикладываться не должна.
\end{itemize}

Решается эта задача (как минимум решалась в самых первых версиях Box2D) с помощью алгоритма Projected Gauss Seidel.
Это обычный Гаусс-Зейдель для СЛАУ, только на каждом шаге дополнительно переменные загоняются обратно в ограничения,
если они были нарушены.

\paragraph{Относительная скорость точки}
Два тела касаются друг друга в одной точке. Эта точка движется как точка первого тела и как точка второго тела.
Скорость точки ($c$) как точки $i$-го тела вычисляется с использованием
скорости центра масс ($v_i$), угловой скорости ($\omega_i$) и координаты центра масс ($p_i$):

\[
    v_{c,i} = v_i - \omega_i \times (c - p_i)
\]

Относительная скорость точки как точки $i$-го тела и как точки $j$-го тела будет $v_{rel} = v_{c,i} - v_{c,j}$.
Её нужно спроецировать на нормаль плоскости касания. Прикладывая силу вдоль этой нормали нужно добиться того,
чтобы относительная скорость стала больше либа равна нулю.

\paragraph{Якобиан}
Допустим имеется $N$ тел и $M$ касаний. Для разрешения касания прикладывается какая-то сила (одинаковая по величине двум телам).
Эти силы мы хотим найти, то есть имеется $M$ неизвестных.
Если бы каждое касание было изолированным от остальных, то проблем бы не было, а так
на одно тело могут влиять несколько неизвестных сил. 

Соответственно, при изменении одной силы могут изменить своё состояние сразу несколько ограничений.
Якобиан -- это матрица, которая проецирует эти силы на ограничения,
показывает как и на какие ограничения влияет такая-то сила.

Поэтому для каждой такой силы в Якобиане будет строка:
\[
\begin{pmatrix}
\dots & n_x & n_y & (r_1 \times n) & \dots & -n_x & -n_y & -(r_2 \times n) & \dots
\end{pmatrix}
\]

\begin{itemize}
    \item $n$ -- нормаль плоскости касания.
    \item $r_1$ -- плечо силы для первого тела.
    \item $r_2$ -- плечо силы для второго тела.
\end{itemize}

Под векторным произведением тут понимается взять компоненту $z$ трехмерного векторного произведения.
Возможно знаки попутаны, но суть такая. Вообще эта матрица получается достаточно разреженной,
чтобы применять на ней соответствующие алгоритмы, но как-нибудь в другой раз.

У каждого тела в 2D три степени свободы: $(x, y, \alpha)$. На них эти силы и проецируем.

\paragraph{Матрица обратных масс}

Сразу же заведем матрицу обратных масс, чтобы дальше ей пользоваться при переводе сил в ускорения (в том числе угловые):

\[
M_{inv} =
\begin{pmatrix}
    M_{inv,1} \\
    & \ddots \\
    & & M_{inv,N}
\end{pmatrix}
\]

\[
M_{inv,i} =
\begin{pmatrix}
    m_i & 0 & 0 \\
    0 & m_i & 0 \\
    0 & 0 & I_i
\end{pmatrix}
\]

$m_i$ и $I_i$ -- массы и моменты инерции соответственно.

\paragraph{LCP для непроникновения}
Попробуем определить $A$ и $b$ для $Ax \ge b$ и остальных условий.
Текущая относительная скорость в точке касания -- $v_{rel,k}$.
Относительная скорость вдоль нормали плоскости касания -- $v_{rel_k} \cdot n_k$.
Хотим $\sum_{i} x_i a_{k,i} + v_{rel_k} \cdot n_k \ge 0$, где
$a_{k,i}$ -- это влияние $i$-ой силы на изменение относительной скорости в $k$-ом касании
Таким образом сразу можно определить вектор $b$:

\[
    b =
    \begin{pmatrix}
        -v_{rel,1} \cdot n_1 \\ \vdots \\ -v_{rel,M} \cdot n_M
    \end{pmatrix}
\]

По хорошему надо так, но в программе у нас еще происходит обработка столкновений, причем дискретно.
Из-за этого столкновения регистрируется уже в тот момент, когда одно тело проникает в другое.
Если использовать continuous collision detection (CCD), то такого конечно же не будет, но мне как-то раз в ленте
попалась запись с GDC или похожего мероприятия про CCD и я ничего не понял.

Поэтому простым хаком будет добавить к относительной скорости небольшое смещение для выталкивания одного тела из другого.
Допустим, $d_k$ -- глубина проникновения, тогда смещение можно задать например вот так: 

\[
    bias_k = \frac{0.2 \cdot (d_k - 0.01)}{DT}
\]

\[
    b =
    \begin{pmatrix}
        -v_{rel,1} \cdot n_1 - bias_1 \\ \vdots \\ -v_{rel,M} \cdot n_M - bias_M 
    \end{pmatrix}
\]

Числа взяты из головы, с физикой это ничего общего не имеет.

\paragraph{}
Заметим, что
с помощью матрицы $J$ можно перевести силы из пространства касаний в пространство тел,
с помощью матрицы $M_{inv}$ можно перевести силы в ускорения,
с помощью матрицы $J^T$ можно перевести ускорения в относительные ускорения в точке контакта.

Таким образом, $A = J M_{inv} J^T$. Матрица получится как раз такой, что
$a_{k,i}$ -- это влияние $i$-ой силы на изменение относительной скорости в $k$-ом касании.

\paragraph{LCP для трения}
Теперь силы прикладываются в плоскости касания. В двухмерном случае просто вдоль прямой.
Якобиан рассчитывается также, только вместо нормали использовать касательную.

Определим $A$ и $b$ вот так:
\[
b =
\begin{pmatrix}
    -v_{rel,1} \cdot t_1 \\ \vdots \\ -v_{rel,M} \cdot t_M
\end{pmatrix}
\]

\[
A = J M J^T
\]

Ограничение для $x_k$:

\[
x_k \in [-x_{n,k} \mu, x_{n,k} \mu]
\]

где $x_{n,k}$ -- нормальная реакция, вычисленная в LCP выше. Перепишем PGS для соблюдения этого ограничения.
Вообще это не соответствует строгому определению LCP (в нём мы хотели $x \ge 0$), но PGS такое решает. 

Немного не понимаю этот момент, выглядит как будто особенности реализации PGS??
В этот раз при прикладывании сил ($x > 0$) мы хотим ($Ax + b = 0$).
Но это так-то не всегда возможно, потому что силы по величине ограничены сверху.
Следовательно когда силы трения не хватает, чтобы занулить относительную скорость,
решения для этой LCP не будет. Но PGS отработает свои $n$ итераций и что-то вернёт.
Предполагаю, что это что-то -- это сила трения скольжения.

То есть PGS пытается занулить скорость, силы $x$ превышают допустимый максимум, обрезаются.
Получается сила трения будет действовать в сторону зануления относительной скорости, но с максимально возможной силой?
Это выглядит как трение скольжения. 

\paragraph{В трехмерном случае}
В трехмерном случае ограничение на величину силы трения превращается в конус и задача становится нелинейной.
В статье (см. ссылку вверху) на 24-й странице предлагается приближать конус с помощью многогранника,
а наличие движения регистрировать с помощью дополнительных индикаторных переменных.
Размер системы из-за этого сильно раздувается, но зато задача становится линейной.

\end{document}